% authority: science_draft
\documentclass[11pt]{article}
\usepackage[margin=1in]{geometry}
\usepackage{amsmath,amssymb,amsthm}
\usepackage{longtable}

\newcommand{\GSC}{\mathsf{GSC}}
\newcommand{\Msrc}{M_{\mathrm{src}}}
\newcommand{\Mgsccand}{M_{\mathrm{src}}^{\GSC\mbox{-}\mathrm{cand}}}
\newcommand{\GenCover}{\mathsf{GenCover}_{\mathrm{src}}^{\GSC}}
\newcommand{\Qsrc}{Q_{\mathrm{src}}^{\GSC}}
\newcommand{\Cert}{\mathsf{Cert}}
\newcommand{\Bottomsrc}{\mathrm{Bottom}_{\mathrm{src}}}
\newcommand{\Csrc}{C_{\mathrm{src}}}
\newcommand{\Fsrc}{F_{\mathrm{src}}}
\newcommand{\Tctx}{T}
\newcommand{\TopT}{T_{\mathrm{top}}}
\newcommand{\AtlasT}{T_{\mathrm{atlas}}}
\newcommand{\MetricT}{T_{\mathrm{metric}}}
\newcommand{\ProperT}{T_{\mathrm{proper}}}
\newcommand{\DetectorT}{T_{\mathrm{detector}}}
\newcommand{\BenchmarkT}{T_{\mathrm{benchmark}}}
\newcommand{\ProcessT}{T_{\mathrm{process}}}
\newcommand{\geff}{g_{\mathrm{eff}}}

\newtheorem{definition}{Definition}
\newtheorem{theorem}{Theorem}
\newtheorem{corollary}{Corollary}
\newtheorem{rulebox}{Rule}

\title{No-Target-Import Criterion for \(M_{\mathrm{src}}^{\mathsf{GSC}\mbox{-}\mathrm{cand}}(E)\) Adoption}
\author{Ontology Formalizer Packet RT-20260614-128}
\date{2026-06-24}

\begin{document}
\maketitle

\section{Control Status}

This artifact is the Phase 5 Ontology Formalizer output for
\texttt{RT-20260614-128}. The selected role is
\texttt{ontology-formalizer@0.2.0}. The parent handoff is
\texttt{handoff-0163}. The target derivation milestone is
\texttt{source\_manifold\_m\_src}.

The purpose is to turn no-target-import discipline into a theorem-shaped
criterion for the current \(\Mgsccand(E)\) adoption attempt. This packet does
not adopt \(\Msrc\), does not construct full \(\Msrc\), does not adopt
\(\mathsf{RegSold}_{\mathrm{src}}^{\GSC}\), does not adopt
\(\mathsf{FVR}_{\mathrm{src}}^{\GSC}\), does not edit canonical ontology TeX,
does not define \(\geff\), does not derive matter coupling, does not derive
Einstein equations, does not promote benchmark status, does not issue Gate
Chair closure, does not claim completed derivation, does not claim future
source-extension impossibility, and does not reject the global theory.

The reviewed candidate remains \(\Mgsccand(E)\) as reviewed source-extension
bridge-candidate data. The result of this packet is:
\[
  \texttt{no\_target\_import\_status =
  criterion\_formalized\_pending\_audit}.
\]

\section{Target Context and Source Construction}

\begin{definition}[Target context]
For audit purposes define a target context
\[
  \Tctx =
  (\TopT,\AtlasT,\MetricT,\ProperT,\DetectorT,\BenchmarkT,\ProcessT),
\]
where \(\TopT\) is target topology, \(\AtlasT\) is target smooth atlas,
\(\MetricT\) is target metric or Lorentzian signature, \(\ProperT\) is target
proper-time or clock semantics, \(\DetectorT\) is empirical detector
semantics, \(\BenchmarkT\) is exact-GR benchmark success data, and
\(\ProcessT\) is validator, registry, role, handoff, generated derivative,
local cache, file order, or commit metadata.
\end{definition}

\begin{definition}[Source construction with audit parameter]
Let
\[
  \Csrc(E;\Tctx) := \Mgsccand(E;\Tctx)
\]
denote the candidate adoption construction when written with an explicit audit
slot for forbidden target context. The intended source-only construction is
\[
  \Csrc(E) := \Mgsccand(E).
\]
The audit slot is not an allowed mathematical input. It is a diagnostic handle
for checking whether a proof accidentally depends on target or process data.
\end{definition}

\begin{definition}[Allowed source data]
Let \(\mathcal{S}(E)\) be the source data generated by
\[
  E,\quad \GenCover,\quad \Qsrc(E),\quad
  \{U_i^{\mathrm{src}}\},\quad \{O_{ij}^{\mathrm{src}}\},\quad
  \{\chi_i^{\mathrm{src}}\},\quad \{\tau_{ij}^{\mathrm{src}}\},\quad
  \Cert(E),\quad \Bottomsrc
\]
and by source-side support, overlap, transition-token, certificate, and
bottom-guard operations explicitly licensed by the adoption theorem envelope.
\end{definition}

\begin{definition}[Source-only factorization]
A candidate adoption proof \(\Pi\) for \(\Mgsccand(E)\) satisfies the
source-only factorization condition when every non-bottom premise and every
non-bottom conclusion in \(\Pi\) factors through \(\mathcal{S}(E)\). Equivalently,
there is a source-only construction \(\Fsrc(E)\) such that
\[
  \Csrc(E;\Tctx)=\Fsrc(E)
\]
for every target context \(\Tctx\) in the diagnostic class. If the equality of
constructed objects is too strong for a later formal setting, the required
replacement is proof-step factorization: each non-bottom assertion of \(\Pi\)
must be expressible without mention of \(\Tctx\).
\end{definition}

\section{No-Target-Import Criterion}

\begin{theorem}[No-target-import criterion]
The candidate adoption proof for \(\Mgsccand(E)\) satisfies no-target-import
only if every non-bottom premise and every non-bottom conclusion factors
through source data \(E\), \(\GenCover\), \(\Qsrc(E)\), source supports, source
overlap predicates, source transition tokens, source certificates, and
\(\Bottomsrc\) guards, and remains invariant under arbitrary replacement of
target topology, target atlas, target metric, proper-time semantics, detector
semantics, benchmark success data, registry metadata, validator status, role
identity, handoff status, generated derivatives, local caches, file order, and
commit metadata.
\end{theorem}

\begin{proof}
Assume the source-only factorization condition. Then each non-bottom premise
and conclusion of the candidate proof \(\Pi\) is the image of data in
\(\mathcal{S}(E)\), or is a bottom-guarded failure branch. Let \(\Tctx\) and
\(\Tctx'\) be arbitrary target contexts. Because \(\Fsrc(E)\) does not take
\(\TopT\), \(\AtlasT\), \(\MetricT\), \(\ProperT\), \(\DetectorT\),
\(\BenchmarkT\), or \(\ProcessT\) as arguments, substituting \(\Tctx'\) for
\(\Tctx\) cannot change any factorized non-bottom assertion. Thus
\[
  \Csrc(E;\Tctx)=\Fsrc(E)=\Csrc(E;\Tctx')
\]
whenever object equality is meaningful, and otherwise the proof-step content
of \(\Pi\) is invariant under the same replacement. Therefore any step whose
truth changes under such replacement fails the criterion. If the failed step
uses target topology, target atlas, target coordinate transitions, target
metric, proper time, detector semantics, or benchmark success, record
\texttt{OB-MSRC-TARGET-IMPORT}. If the failed step uses validator, registry,
role, handoff, generated derivative, local cache, file order, or commit status
as proof, record \texttt{OB-MSRC-PROCESS-AUTHORITY-LAUNDERING}. If the failed
step converts finite/local survival into global adoption, retain
\texttt{OB-MSRC-LOCAL-GLOBAL-GAP} as the inference-blocking label.
\end{proof}

\begin{corollary}[Criterion status]
This packet formalizes the no-target-import criterion for the current adoption
attempt. It does not prove that a future integrated adoption theorem satisfies
the criterion, and it does not adopt \(\Msrc\).
\end{corollary}

\section{Proof Obligations by Import Class}

\begin{longtable}{|p{0.22\linewidth}|p{0.47\linewidth}|p{0.20\linewidth}|}
\hline
\textbf{Import class} & \textbf{Required demonstration} & \textbf{Failure label} \\
\hline
Target topology & No target open set, target neighborhood, or target manifold topology appears in premises, definitions, or proof steps. & \texttt{OB-MSRC-TARGET-IMPORT} \\
\hline
Target atlas & \(\chi_i^{\mathrm{src}}\) are source support names or tokens only, not adopted target coordinate charts. & \texttt{OB-MSRC-TARGET-IMPORT} \\
\hline
Target transition maps & \(\tau_{ij}^{\mathrm{src}}\) are source transition tokens, not inherited smooth transition maps. & \texttt{OB-MSRC-TARGET-IMPORT} \\
\hline
Target metric & No metric tensor, distance, line element, signature, causal cone, or proper-time object is used. & \texttt{OB-MSRC-TARGET-IMPORT} \\
\hline
Empirical detector semantics & No detector response, observed clock behavior, or empirical observer protocol is used as a proof input. & \texttt{OB-MSRC-TARGET-IMPORT} \\
\hline
Benchmark success & No exact-GR recovery assumption or benchmark success datum is used as source-side evidence. & \texttt{OB-MSRC-TARGET-IMPORT} \\
\hline
Process authority & Validation, registry, role, handoff, generated derivative, local cache, file order, and commit status do not enter the mathematical proof. & \texttt{OB-MSRC-PROCESS-AUTHORITY-LAUNDERING} \\
\hline
Local witness overread & Finite/local witness status is not treated as a global theorem or adoption proof. & \texttt{OB-MSRC-LOCAL-GLOBAL-GAP} \\
\hline
\end{longtable}

\section{Checklist Interpretation}

The no-target-import guard map remains a checklist and routing surface. This
artifact supplies the mathematical criterion that the checklist should test
against when a later integrated theorem attempt exists. A checklist pass is
not by itself a theorem. A validator pass is not by itself a theorem. A
registry row is not by itself a theorem.

\begin{rulebox}[Fail-closed target-context branch]
If any non-bottom proof step requires an element of \(\Tctx\), the adoption
route does not silently continue. It must either return \(\Bottomsrc\), record
\texttt{OB-MSRC-TARGET-IMPORT}, record
\texttt{OB-MSRC-PROCESS-AUTHORITY-LAUNDERING}, or preserve
\texttt{OB-MSRC-LOCAL-GLOBAL-GAP} for finite/local overread. No such failure
authorizes \(\Msrc\) adoption or downstream GR claims.
\end{rulebox}

\section{Relationship to Prior Phases}

\begin{longtable}{|p{0.19\linewidth}|p{0.34\linewidth}|p{0.34\linewidth}|}
\hline
\textbf{Prior phase} & \textbf{Relevant result} & \textbf{Phase 5 use} \\
\hline
Phase 1 & Minimum theorem envelope A1--A9. & Supplies the adoption obligations whose proof steps must factor through source data. \\
\hline
Phase 2 & \texttt{site\_coverage} selected as source topology analogue route. & Must remain a source-only route and not import target open sets. \\
\hline
Phase 3 & Source transition-token groupoid candidate conditions. & \(\tau_{ij}^{\mathrm{src}}\) must remain symbolic source tokens rather than target smooth maps. \\
\hline
Phase 4 & Finite/local survival does not imply global adoption. & Blocks finite/local witness overread through \texttt{OB-MSRC-LOCAL-GLOBAL-GAP}. \\
\hline
\end{longtable}

\section{Acceptance Criterion}

The Phase 5 acceptance criterion is satisfied because the packet defines:

\begin{itemize}
  \item the target context \(\Tctx\);
  \item the diagnostic construction \(\Csrc(E;\Tctx)\);
  \item source-only factorization through \(\mathcal{S}(E)\);
  \item the no-target-import criterion theorem;
  \item import-class proof obligations; and
  \item failure labels for target import, process-authority laundering, and
        finite/local overread.
\end{itemize}

The packet's status is
\[
  \texttt{criterion\_formalized\_pending\_audit}.
\]

\section{Next Route}

The next legal packet under the implementation plan is Phase 6: one bounded
finite checker or model-search packet used as a counterexample factory for the
source-only adoption theorem obligations. It must not treat checker success as
proof and must not route to \(\geff\), matter coupling, Einstein equations,
benchmark promotion, Gate Chair closure, completed derivation, future
source-extension impossibility, or global theory rejection.

The criterion itself should be audited by a Smuggling Auditor after Phase 7
integrates it into an adoption theorem attempt, unless a later Director
explicitly selects the criterion as the active audited candidate.

\section{Conclusion}

No-target-import is now represented as an auditable source-only factorization
criterion. This narrows the adoption theorem burden but does not discharge it.
Full \(\Msrc\) adoption remains not granted. The same-milestone continuation
remains open.

\begin{thebibliography}{9}

\bibitem{plan}
The AEther-Flow Research Project. (2026, June 24). \emph{Recommendations
implementation plan for continue-research v4} [Internal implementation plan].

\bibitem{guard-map}
The AEther-Flow Research Project. (2026, June 23). \emph{No-target-import
guard map} [Internal control note].

\bibitem{theorem-envelope}
The AEther-Flow Research Project. (2026, June 24). \emph{Minimum source-only
adoption theorem envelope for M\_src GSC candidate} [Internal research-control
artifact].

\bibitem{transition-conditions}
The AEther-Flow Research Project. (2026, June 24). \emph{Source transition-token
groupoid conditions for M\_src GSC candidate} [Internal research-control
artifact].

\bibitem{finite-local}
The AEther-Flow Research Project. (2026, June 24). \emph{Finite/local witness
classification and global M\_src adoption gap} [Internal research-control
artifact].

\end{thebibliography}

\end{document}
